%% Cross-Species Alignment via Cycle-Consistent Adapter NMF
%% Migrated from singlet-intelligence/docs/architecture/cross_species.md
\documentclass[10pt,twocolumn,letterpaper]{article}
\usepackage{../../templates/singletai-preprint}
\usepackage{graphicx}
\usepackage{amsmath,amssymb}
\usepackage{booktabs}
\usepackage{algorithm}
\usepackage{algpseudocode}
\usepackage{natbib}
\bibliographystyle{unsrtnat}

% ── Metadata ─────────────────────────────────────────────────────
\title{Cross-Species Alignment via Cycle-Consistent Adapter NMF}
\author{
  Zach DeBruine\textsuperscript{1,*}
  \\[6pt]
  \textsuperscript{1}Department of Computing, Grand Valley State University, Allendale, MI 49401
}

\begin{document}
\maketitle
\correspondingauthor{debruinz@gvsu.edu}

\begin{abstract}
We present a cycle-consistent adapter framework for aligning non-negative
matrix factorization (NMF) models across species without paired observations,
gene homology tables, or cell type labels. Given independently pre-trained NMF
models for two species, we learn a non-negative adapter matrix $A$ that maps
program activities between species while preserving reconstruction fidelity.
The transpose parameterization $B = A^T$ provides a strong inductive bias
toward permutation-like mappings, halving parameters and guaranteeing symmetric
alignment. We derive multiplicative update rules for all variables---adapter,
activities, and gene dictionaries---ensuring non-negativity by construction.
Three algorithms are presented: one-sided alignment with a frozen reference
(Algorithm~1), joint training (Algorithm~2), and a mini-batch extension
(Algorithm~3) requiring only ${\sim}1.5$~GB memory regardless of dataset size.
The framework extends naturally to $N > 2$ species via a hub-and-spoke
architecture with human as the reference frame.
\end{abstract}

% ── Section 1 ────────────────────────────────────────────────────
\section{Problem Statement}

We have independently pre-trained NMF models for two species (human, mouse).
Each model decomposes its atlas into gene programs:
\begin{equation}
  X_h \approx W_h H_h, \qquad X_m \approx W_m H_m
\end{equation}
where $W_s \in \mathbb{R}_{\geq 0}^{g_s \times k_s}$ are gene program
dictionaries and $H_s \in \mathbb{R}_{\geq 0}^{k_s \times n_s}$ are program
activities.

\textbf{The problem:} Learn a mapping between program spaces such that
biologically equivalent programs are aligned. We never observe paired
mouse--human cells. We have no gene homology tables. We have no cell type
labels as training signal. Only the internal geometry of each program space.

\textbf{Key constraints:}
\begin{itemize}
  \item All matrices remain non-negative throughout (NMF structure preserved).
  \item The human model is the \textbf{reference frame}---its program space
        defines the universal vocabulary.
  \item The mouse model may be fine-tuned to improve alignment while
        maintaining reconstruction quality.
  \item The mapping must be learnable via multiplicative updates (no
        gradient-based optimizers, no negative intermediate values).
\end{itemize}

% ── Section 2 ────────────────────────────────────────────────────
\section{Notation and Setup}

\begin{table}[h]
\centering\small
\begin{tabular}{@{}lll@{}}
\toprule
Symbol & Shape & Description \\
\midrule
$X_s$ & $g_s \times n_s$ & Expression matrix for species $s$ \\
$W_s$ & $g_s \times k_s$ & Gene program dictionary \\
$H_s$ & $k_s \times n_s$ & Program activities \\
$A$ & $k_h \times k_m$ & Adapter: mouse $\to$ human \\
$G_s = H_s H_s^T$ & $k_s \times k_s$ & Program co-activation Gram matrix \\
$C_s = G_s / n_s$ & $k_s \times k_s$ & Program covariance \\
$S = A^T A$ & $k_m \times k_m$ & Mouse-side cycle operator \\
$P = A A^T$ & $k_h \times k_h$ & Human-side cycle operator \\
\bottomrule
\end{tabular}
\caption{Notation. Typical dimensions: $g \approx 30{,}000$ genes,
$k \approx 1{,}024$--$2{,}048$ programs, $n \approx 10^7$--$10^9$ cells.}
\end{table}

% ── Section 3 ────────────────────────────────────────────────────
\section{Should Adapters Be Transposes?}

\subsection{The Question}

The previous formulation used two independent adapters:
$A \in \mathbb{R}_{\geq 0}^{k_h \times k_m}$ (mouse$\to$human) and
$B \in \mathbb{R}_{\geq 0}^{k_m \times k_h}$ (human$\to$mouse). Should we
instead set $B = A^T$?

\subsection{Analysis}

\textbf{Theorem (Berman--Plemmons).} A non-negative matrix with a
non-negative inverse is a \emph{monomial matrix}---a product of a permutation
matrix and a positive diagonal matrix.

If we require both $BA \approx I_{k_m}$ and $AB \approx I_{k_h}$ with
$A, B \geq 0$, then $A$ is constrained to be approximately monomial. With
$B = A^T$:
\begin{itemize}
  \item $A^T A \approx I_{k_m}$: columns of $A$ are approximately orthonormal.
  \item $A A^T \approx I_{k_h}$: rows of $A$ are approximately orthonormal.
  \item Combined with non-negativity: $A$ is approximately a
        \textbf{permutation matrix} (or a padded permutation if
        $k_h \neq k_m$).
\end{itemize}

A non-negative approximately-orthogonal matrix is approximately a permutation.
This is a \emph{feature}, not a bug---it provides an extremely strong inductive
bias toward crisp 1-to-1 program correspondences, which is the biologically
expected structure for well-conserved programs.

\subsection{Many-to-Many Mappings}

If mouse program~12 genuinely corresponds to a blend of human programs
$\{7, 23\}$, the transpose parameterization cannot represent this as a clean
mapping. However:

\begin{enumerate}
  \item \textbf{Fine-tuning $W_m$ absorbs this.} The $W_m$ update can
        reorganize mouse programs so that two distinct programs each correspond
        to one human program, while reconstruction loss ensures fidelity.
  \item \textbf{The cycle residual diagnoses it.} Programs that cannot be
        cleanly aligned produce high cycle reconstruction error, flagging them
        as species-divergent---valuable biological signal.
  \item \textbf{The rectangular case handles dimensionality mismatch.} If
        $k_h > k_m$, $A$ maps each mouse program to a direction in
        $\mathbb{R}^{k_h}$. The constraint $A^T A \approx I_{k_m}$ says
        these directions are orthonormal. The residual
        $\|(I - AA^T) H_h\|_F^2$ identifies human programs with no mouse
        counterpart.
\end{enumerate}

\subsection{Recommendation}

\textbf{Use $B = A^T$ as the default parameterization.}

\begin{table}[h]
\centering\small
\begin{tabular}{@{}lll@{}}
\toprule
Property & Independent $A, B$ & Transpose $B = A^T$ \\
\midrule
Parameters & $2 k_h k_m$ & $k_h k_m$ \\
Inductive bias & Weak & Strong (permutation) \\
Identifiability & Degenerate solutions & Well-constrained \\
Symmetry & May disagree & Guaranteed \\
Divergent programs & Absorbed & Flagged by residual \\
\bottomrule
\end{tabular}
\end{table}

% ── Section 4 ────────────────────────────────────────────────────
\section{Algorithm 1: One-Sided Alignment}

\subsection{Setting}

The human model $(W_h, H_h)$ is frozen. We learn: $A$ (adapter),
$H_m$ (mouse activities, fine-tuned under cycle pressure), and optionally
$W_m$ (mouse dictionary, anchored to pretrained $W_m^{(0)}$).

\subsection{Objective}

\begin{align}
  \mathcal{L} &= \underbrace{\|X_m - W_m H_m\|_F^2}_{\text{reconstruction}}
  + \lambda_{\text{cyc}} \Big(
    \underbrace{\|H_m - A^T\!A\,H_m\|_F^2}_{\text{mouse cycle}}
    + \underbrace{\|H_h - AA^T\!H_h\|_F^2}_{\text{human cycle}}
  \Big) \nonumber\\
  &\quad + \lambda_{\text{str}}
    \underbrace{\|C_h - A\,C_m\,A^T\|_F^2}_{\text{structural alignment}}
  + \mu \underbrace{\|W_m - W_m^{(0)}\|_F^2}_{\text{anchor}}
\end{align}

\textbf{Loss terms:}
\begin{itemize}
  \item \textbf{Reconstruction:} Mouse model must still explain its data.
  \item \textbf{Mouse cycle} ($A^T A \approx I$): Round-trip
        mouse$\to$human$\to$mouse recovers original activities.
  \item \textbf{Human cycle} ($AA^T \approx I$): Round-trip
        human$\to$mouse$\to$human recovers activities. For $k_m < k_h$ this
        cannot go to zero---the residual identifies human-specific programs.
  \item \textbf{Structural alignment:} Co-activation geometry of mouse
        programs, when mapped to human space, should match human geometry.
  \item \textbf{Anchor:} Optional proximity penalty keeping $W_m$ near
        pretrained value. Set $\mu = 0$ for unconstrained fine-tuning.
\end{itemize}

\subsection{Multiplicative Update Rules}

Each update is derived by splitting the gradient $\partial \mathcal{L} / \partial \theta$
into non-negative and non-positive components, yielding the Lee--Seung form:
\begin{equation}
  \theta \leftarrow \theta \odot
    \frac{[\nabla_\theta \mathcal{L}]^-}{[\nabla_\theta \mathcal{L}]^+}
\end{equation}

\paragraph{$H_m$ Update.}
\begin{equation}
  \boxed{H_m \leftarrow H_m \odot
    \frac{W_m^T X_m + 2\lambda_{\text{cyc}}\,S\,H_m}
         {W_m^T W_m\,H_m + \lambda_{\text{cyc}}(I + S^2)\,H_m}}
\end{equation}
where $S = A^T A$ is $k_m \times k_m$ (precomputed once per iteration).

\paragraph{$W_m$ Update.}
\begin{equation}
  \boxed{W_m \leftarrow W_m \odot
    \frac{X_m H_m^T + \mu\,W_m^{(0)}}{W_m\,G_m + \mu\,W_m}}
\end{equation}
where $G_m = H_m H_m^T$. Standard NMF update with proximity anchor.

\paragraph{$A$ Update.}
Three gradient contributions (mouse cycle, human cycle, structural alignment)
combine to:
\begin{equation}
  \boxed{A \leftarrow A \odot
    \frac{2\lambda_{\text{cyc}}(AG_m + G_h A) + 2\lambda_{\text{str}}\,C_h A C_m}
         {\lambda_{\text{cyc}}\big(A(SG_m\!+\!G_m S) + (PG_h\!+\!G_h P)A\big)
          + 2\lambda_{\text{str}}\,(AC_m A^T)\,A\,C_m}}
\end{equation}
where $S = A^T A$, $P = A A^T$, $G_s = H_s H_s^T$, $C_s = G_s / n_s$
(all precomputed). All terms are non-negative.

\subsection{Algorithm Pseudocode}

\begin{algorithm}[h]
\caption{One-Sided Cycle-Consistent Adapter NMF}
\begin{algorithmic}[1]
\Require $X_m$, frozen $(W_h, H_h)$, pretrained $(W_m^{(0)}, H_m^{(0)})$,
         $\lambda_{\text{cyc}}, \lambda_{\text{str}}, \mu, T$
\State $W_m \gets W_m^{(0)}$; $H_m \gets H_m^{(0)}$;
       $A \gets \text{random non-neg, col-normalized}$
\State Precompute $G_h \gets H_h H_h^T$; $C_h \gets G_h / n_h$
\For{$t = 1, \ldots, T$}
  \State $G_m \gets H_m H_m^T$; $S \gets A^T A$
  \State Update $H_m$ (Eq.~3)
  \State $G_m \gets H_m H_m^T$ \Comment{recompute}
  \State Update $W_m$ (Eq.~4); normalize columns, scale $H_m$
  \State Recompute $G_m, C_m, S, P$
  \State Update $A$ (Eq.~5)
  \State Log losses: $L_{\text{recon}}, L_{\text{cyc\_m}}, L_{\text{cyc\_h}}, L_{\text{str}}$
\EndFor
\State \Return $W_m, H_m, A$
\end{algorithmic}
\end{algorithm}

\subsection{Inference}

Given a new mouse cell $x_m$:

\textbf{Step~1} --- Project to mouse programs:
$h_m = \arg\min_{h \geq 0} \|x_m - W_m h\|^2$.

\textbf{Step~2} --- Map to human program space:
$h_h^{\text{pred}} = A\,h_m$.

\textbf{Step~3} --- Decode through human dictionary:
$\hat{x}_h = W_h\,h_h^{\text{pred}}$.

\textbf{Confidence per program:} $(A^T A)_{jj} \approx 1$ means program $j$
maps cleanly; $(A^T A)_{jj} \ll 1$ means species-divergent.
$(I - AA^T)_{ii} \approx 1$ identifies human programs absent from the mouse
mapping.

% ── Section 5 ────────────────────────────────────────────────────
\section{Algorithm 2: Joint Training}

\subsection{Motivation}

In Algorithm~1, the human model is frozen. Joint training allows both species'
program spaces to co-adapt, potentially finding a geometry that is
simultaneously faithful to each species' data and mutually alignable.

\subsection{Objective}

\begin{align}
  \mathcal{L} &= \|X_h - W_h H_h\|_F^2 + \|X_m - W_m H_m\|_F^2 \nonumber\\
  &\quad + \lambda_{\text{cyc}} \big(
    \|H_m - S\,H_m\|_F^2 + \|H_h - P\,H_h\|_F^2
  \big) \nonumber\\
  &\quad + \lambda_{\text{str}} \big(
    \|C_h - A\,C_m\,A^T\|_F^2
    + \|C_m - A^T C_h\,A\|_F^2
  \big)
\end{align}

The new updates for the human model are:
\begin{equation}
  H_h \leftarrow H_h \odot
    \frac{W_h^T X_h + 2\lambda_{\text{cyc}}\,P\,H_h}
         {W_h^T W_h\,H_h + \lambda_{\text{cyc}}(I + P^2)\,H_h}
\end{equation}
\begin{equation}
  W_h \leftarrow W_h \odot \frac{X_h H_h^T}{W_h\,G_h}
\end{equation}

The joint $A$ update incorporates the second structural term, yielding
identical numerator terms ($C_h A C_m$ appears twice, simplifying to
$4\lambda_{\text{str}} C_h A C_m$).

\subsection{Risk: Co-Adaptation Collapse}

Joint training risks a degenerate solution where both models reorganize
programs so alignment is trivial ($A \to I$) but programs lose biological
meaning.

\textbf{Why reconstruction prevents this:} The reconstruction losses are
computed in \textbf{gene space} ($g \times n$), not program space. With
$g \approx 30{,}000$ genes and $n \approx 10^8$ cells, reconstruction
dominates by orders of magnitude.

\textbf{Annealing is critical:} Start with $\lambda_{\text{cyc}} =
\lambda_{\text{str}} = 0$ and linearly ramp up. Final $\lambda$ values
should contribute $\leq 10\%$ of total loss.

\textbf{Recommendation:} Run Algorithm~1 first. Use Algorithm~2 only if
cycle residual on conserved programs is too high.

% ── Section 6 ────────────────────────────────────────────────────
\section{Computational Analysis}

\subsection{The $C$ Matrix Is Not Large}

All adapter operations are in \textbf{program space}, not gene or cell space:

\begin{table}[h]
\centering\small
\begin{tabular}{@{}llr@{}}
\toprule
Matrix & Dimensions & Memory \\
\midrule
$C_s, G_s$ & $k \times k$ & ${\sim}33$~MB \\
$A$ & $k_h \times k_m$ & ${\sim}33$~MB \\
$S, P$ & $k \times k$ & ${\sim}33$~MB each \\
\bottomrule
\end{tabular}
\caption{Memory for $k = 2048$, float64. All fit in L2 cache.}
\end{table}

\subsection{Computing $G = HH^T$ at Scale}

$G_m = \sum_{i=1}^{n_m} h_i h_i^T$ via streaming accumulation. Memory:
$O(k^2)$. For $k = 2048$, $n = 10^8$: ${\sim}100$~seconds at 4~TFLOPS
(single A100)---negligible compared to NMF reconstruction updates.

\subsection{Actual Bottlenecks}

\begin{table}[h]
\centering\small
\begin{tabular}{@{}llr@{}}
\toprule
Operation & FLOPs & Note \\
\midrule
$W_m^T X_m$ & $O(k \cdot \text{nnz})$ & Sparse \\
$W_m^T W_m H_m$ & $O(k^2 n)$ & Streamed \\
Adapter update & $O(k^3)$ & Negligible \\
\bottomrule
\end{tabular}
\caption{Adapter adds ${\sim}10^{10}$ FLOPs per iteration vs.\
${\sim}10^{13}$ for NMF---three orders of magnitude smaller.}
\end{table}

% ── Section 7 ────────────────────────────────────────────────────
\section{Mini-Batch Extension}

To avoid storing the full $H_m$ ($k \times n$ can exceed 1~TB), we adopt an
epoch-based mini-batch algorithm that accumulates sufficient statistics.

\begin{algorithm}[h]
\caption{Mini-Batch Cycle-Consistent Adapter NMF}
\begin{algorithmic}[1]
\For{each epoch}
  \State $XH \gets 0 \in \mathbb{R}^{g_m \times k_m}$;
         $G_m \gets 0 \in \mathbb{R}^{k_m \times k_m}$
  \For{each batch $B$ of cells ($|B| \sim 10^3$--$10^4$)}
    \State Load $X_B = X_m[:, B]$ from disk
    \State $H_B \gets \text{NNLS}(W_m, X_B)$
    \State Apply MU update to $H_B$ (Eq.~3)
    \State $XH \mathrel{+}= X_B H_B^T$; \;
           $G_m \mathrel{+}= H_B H_B^T$
  \EndFor
  \State Update $W_m$ using accumulated $XH, G_m$
  \State Update $A$ using accumulated $G_m, C_m$
\EndFor
\end{algorithmic}
\end{algorithm}

\textbf{Memory budget:} ${\sim}1.5$~GB regardless of dataset size. The full
$H_m$ is never materialized.

\begin{table}[h]
\centering\small
\begin{tabular}{@{}llr@{}}
\toprule
Object & Size & Lifetime \\
\midrule
$X_B$ (batch) & ${\sim}300$~MB & Per-batch \\
$H_B$ (batch) & ${\sim}20$~MB & Per-batch \\
$XH$ accumulator & ${\sim}500$~MB & Per-epoch \\
$W_m, A$ & ${\sim}530$~MB & Persistent \\
\bottomrule
\end{tabular}
\end{table}

% ── Section 8 ────────────────────────────────────────────────────
\section{Initialization}

\textbf{$W_m, H_m$:} Start from pretrained mouse NMF (warm start).

\textbf{$A$---Option~A (Random):} $A_{ij} \sim \text{Uniform}(0, 1/\sqrt{k_m})$,
columns normalized to unit $L_2$.

\textbf{$A$---Option~B (Spectral matching):} Compute eigendecompositions
$C_s = U_s \Lambda_s U_s^T$, initialize
$A \approx U_h[:,1{:}k_m] \cdot U_m[:,1{:}k_m]^T$, rectify to non-negative.
Use when $k > 512$.

\textbf{$A$---Option~C (Diagonal):} $A = \alpha I$ if $k_h = k_m$. Fast
convergence if programs happen to be index-aligned.

% ── Section 9 ────────────────────────────────────────────────────
\section{Convergence and Normalization}

\textbf{Monotonic decrease:} Each MU is derived via auxiliary functions
(Lee \& Seung, 2001). Block-coordinate structure guarantees convergence to a
stationary point.

\textbf{Column normalization:} Normalize $W_m$ columns to unit $L_2$ after
each update with compensating $H_m$ rescaling. Do \textbf{not} normalize $A$
---the cycle loss regulates its scale.

\textbf{Numerical stability:} Add $\epsilon = 10^{-10}$ to both numerator and
denominator. Rarely triggered with warm starts.

% ── Section 10 ───────────────────────────────────────────────────
\section{Validation}

\subsection{Intrinsic Metrics (No Labels)}

\begin{table}[h]
\centering\small
\begin{tabular}{@{}llr@{}}
\toprule
Metric & Target \\
\midrule
Reconstruction error & $\leq 1.05\times$ pretrained \\
Mouse cycle fidelity & $< 0.05$ \\
Human cycle fidelity & Decreasing \\
Structural alignment & $< 0.1$ \\
Adapter orthogonality & $< 0.1$ \\
Sparsity of $A$ (Gini) & $> 0.8$ \\
\bottomrule
\end{tabular}
\end{table}

\subsection{Extrinsic Metrics (Held-Out Labels)}

Cell type annotations provide post-hoc validation:
\begin{itemize}
  \item \textbf{Transfer accuracy:} Mean $h_m$ per mouse cell type, map via
        $A$, match to nearest human centroid.
  \item \textbf{Program coherence:} Gene set enrichment overlap for aligned
        program pairs using ortholog mappings.
  \item \textbf{Bidirectional consistency:} Transfer accuracy should be
        similar in both directions.
\end{itemize}

\subsection{Biological Sanity Checks}

Oxidative phosphorylation, ribosomal, and cell cycle programs should align
with high confidence (deeply conserved). Immune programs should show partial
alignment with species-specific residuals. Cycle residual magnitude should
correlate with evolutionary divergence.

% ── Section 11 ───────────────────────────────────────────────────
\section{Summary of Design Decisions}

\begin{table}[h]
\centering\small
\begin{tabular}{@{}lll@{}}
\toprule
Decision & Choice & Rationale \\
\midrule
Parameterization & $B = A^T$ & Halves params; permutation bias \\
Human model & Frozen or joint & Frozen is safer \\
Loss terms & Recon+cycle+struct & Distinct roles \\
Optimization & Multiplicative & Preserves non-negativity \\
$C$ matrix & $k \times k$ & 33~MB regardless of $n$ \\
Scalability & Mini-batch & 1.5~GB memory \\
Annealing & $\lambda$ ramp from 0 & Recon first, then align \\
Initialization & Pretrained + spectral & Warm start \\
\bottomrule
\end{tabular}
\end{table}

% ── Appendix ─────────────────────────────────────────────────────
\appendix
\section{Extension to $N > 2$ Species}

The framework extends to $N$ species by designating human as the hub:
$A^{(s)} \in \mathbb{R}_{\geq 0}^{k_h \times k_s}$ for each species
$s \neq \text{human}$. Cross-species transfer routes through human:
\begin{equation}
  \hat{h}_{\text{zf}} = (A^{(\text{zf})})^T \cdot A^{(\text{mouse})} \cdot h_{\text{mouse}}
\end{equation}

This hub-and-spoke architecture avoids $O(N^2)$ pairwise adapters, keeps human
as the interpretable reference frame, and allows adding new species without
retraining existing adapters.

\section{Gradient Derivations}

\subsection{Mouse Cycle Term}

$f = \|(I - A^T A) H_m\|_F^2 = \text{tr}\big((I - S)^2 G_m\big)$.

Expanding: $f = \text{tr}(G_m) - 2\,\text{tr}(SG_m) + \text{tr}(S^2 G_m)$.

Using $\frac{\partial}{\partial A}\text{tr}(A^T A\,G_m) = 2AG_m$ and
$\frac{\partial}{\partial A}\text{tr}(S^2 G_m) = 2A(SG_m + G_m S)$:
\begin{equation}
  \nabla_A f = -4AG_m + 2A(SG_m + G_m S)
\end{equation}

\subsection{Human Cycle Term}

$f = \|(I - AA^T) H_h\|_F^2 = \text{tr}\big((I - P)^2 G_h\big)$.

Using $\frac{\partial}{\partial A}\text{tr}(AA^T G_h) = 2G_h A$ and
$\frac{\partial}{\partial A}\text{tr}(P^2 G_h) = 2(PG_h + G_h P)A$:
\begin{equation}
  \nabla_A f = -4G_h A + 2(PG_h + G_h P)A
\end{equation}

\subsection{Structural Alignment Term}

$f = \|C_h - AC_m A^T\|_F^2$. Using the standard identity:
\begin{equation}
  \nabla_A f = -4C_h A C_m + 4(AC_m A^T)\,A\,C_m
\end{equation}

\bibliography{../../shared/refs}
\end{document}
